\section{评估设计}
\label{sec:evaluation}
我们通过四个问题评估已实现的边界：(RQ1) 选定的契约测试与运行控制检查是否保持了其声明的边界；(RQ2) 有界上下文是否交付了标注要求的源码；(RQ3) 较弱谓词会掩盖哪些就绪判定错误；(RQ4) 响应大小、可见决策记录和延迟如何随所测试的预算变化？研究单位是一次运行时调用，不包含语言模型调用。

\subsection{语料与执行协议}
诊断包含开发者编写的 60 个合成场景：Rust 27 个，Go、TypeScript 和 Python 各 11 个。测试夹具的文件数介于 26 和 644 之间。每个场景均显式标注了 Gold 源码要求，包括带路径限定的身份与原始源码范围。在 59 个可回答场景中，必需原始源码的并集很小：每个场景为 61--220 字节，中位数为 162 字节。一个场景不可回答，另一个场景可回答但只读，其余 58 个场景符合编辑条件。这些性质使语料适合契约诊断，但源码交付的召回任务较为简单。

我们使用预算 $B\in\{1000,2000,4000\}$、两个执行框架缓存阶段，并在每个场景和条件下重复七次。由此得到 $60\times3\times2\times7=2520$ 次实测查询，另有 180 次独立预热查询。“冷态”创建新的工具执行框架；“预热态”使用已预热的状态。实验不清空操作系统缓存。预算顺序固定，缓存阶段的执行顺序按场景交替。重复运行刻画的是同一批任务的执行变异，不会增加独立任务的数量。

主运行使用提交 \code{584abd88} 对应的源码归档、软件包版本 0.9.6，以及保留的 debug 可执行文件。运行主机为 Apple M5，物理核和逻辑核均为十个，内存为 16\,GiB，操作系统为 arm64 上的 macOS 27.2。运行时的工具链探测报告为 rustc 1.99.1；它不是对编译工具链的独立证明。执行框架配置为四个槽位。主机负载未受控制。完整的源码与可执行文件身份、协议和恢复历史记录于附录~\ref{sec:artifact} 以及随附证据清单中。

\subsection{交付与就绪指标}
在每种条件下，420 次实测尝试包含 686 个必需源码实例、413 次可回答观察，以及 406 次符合编辑条件的观察。我们针对最终响应评估以下嵌套的交付谓词：
\begin{description}[style=unboxed,leftmargin=0pt,font=\bfseries]
\item[身份。] 响应包含标注的、带路径限定的目标。
\item[原始函数体。] 必需源码未经截断地交付；该谓词本身已经检查源码窗口哈希。
\item[文件元数据。] 关联的 \code{files[]} 条目还报告当前 SHA-256。
\item[编辑输入。] 响应满足完整的编辑输入契约，包括写入资格。
\end{description}
前三项以源码实例为单位的指标，在每种条件下均以 686 项要求为分母；编辑输入成功率则以 406 次符合条件的观察为分母。另有一次按尝试统计的就绪分析，在 413 次可回答观察中比较被判为阳性的次数。明确保留这些分母，能够避免将源码召回率误报为编辑就绪率。

我们还测量最终响应中是否包含序列化的 \code{DecisionBatch}，即一种可选的结构化决策记录。其分母为全部 420 次尝试。暴露率衡量调用方能够检查到什么，并不代表决策正确性。原始报告还保留了影子上下文充分性评分；这是一项独立诊断，不改变暴露率的定义。

\subsection{控制检查、消融与计时}
十一项运行控制检查覆盖四种语言各自的语法图、受保护编辑与失效处理、只读准入、基于相同哈希的竞争写入、工作区边界、不完整图的报告、语法缓存压力，以及工作区相对路径下的验证请求合并（singleflight）。它们与主试验一同考察选定机制。

我们还在同一保留可执行文件上，按确切测试名分别重放 22 项已有的选定反例测试，每项执行一次。它们覆盖证据身份与权威性（11 项）、生产者冲突（2 项）、依赖归属（2 项）、范围与可观测性（2 项）以及评分（5 项）。每次调用都记录参数、结果，以及运行前后的可执行文件哈希。重放前后还核对源码清单与主原始报告。先前使用另一构建的执行仍作为历史记录保留；本文报告的结果来自这次新重放。

谓词分析重新组合了对最终响应评分所得并保留下来的逐样本字段；完整的序列化工具响应未被归档。这是对就绪定义进行的观察性消融，并非禁用某个模块的运行时实验。开销方面，我们报告上下文调用的序列化字节数与墙钟延迟。计时不包含编译、夹具构造、评分、模型推理、网络传输或客户端重试。表~\ref{tab:cost} 汇总每种条件下的全部 420 次观察。图~\ref{fig:latency} 则绘制 60 个场景中位数，每个中位数由七次重复得到，从而避免在视觉呈现中把反复执行的夹具当作独立任务。本文报告描述性分布，不给出推断性置信区间。
